\subsubsection{Propuesta de solución [Pendientes: 1. concretar el sector y la cobertura de escenarios, 2. definir selección y agregación, 3. fijar horizonte de consulta y servicio meteorológico]}\label{sec:solucion}
Se propone desarrollar una aplicación de escritorio sustentada en una biblioteca de escenarios precalculados con WFDS. Su preparación comprenderá la configuración del simulador, la elaboración de entradas territoriales compatibles y la ejecución anticipada de casos mediante recursos de cómputo de alto desempeño, de acuerdo con las capacidades que se comprueben. Las igniciones de referencia se distribuirán con densidad variable guiada por amenaza de ignición, de manera que las zonas de mayor propensión tengan más puntos y las de menor propensión conserven cobertura con puntos más separados. Esta distribución corresponde a una decisión del proyecto, cuyo antecedente territorial es el trabajo de Morales Salguero~\cite{morales2017}, y debe distinguirse de la discretización tridimensional utilizada para resolver cada simulación. Los resultados se almacenarán con su ubicación, condiciones de entrada, tiempos y procedencia, de modo que la aplicación pueda recuperar escenarios sin ejecutar de nuevo el cálculo físico completo durante una consulta.

En el uso previsto, el usuario proporcionará una ubicación inicial o un perímetro observado y la aplicación incorporará condiciones meteorológicas actuales obtenidas mediante una API. A partir de esos datos se seleccionarán escenarios próximos espacialmente y compatibles con las condiciones representadas en la biblioteca. Cuando exista un perímetro, la consulta considerará igniciones cercanas al borde activo y agregará sus resultados espaciales, excluyendo el interior que ya se identifica como quemado. La salida se mostrará como una animación georreferenciada para un horizonte limitado, denominado $\Delta t$, y podrá actualizarse ante nuevas observaciones o cambios relevantes del viento. La forma de obtener y verificar el perímetro, las reglas de semejanza, las ponderaciones y la estabilidad temporal que determine $\Delta t$ deberán precisarse y comprobarse durante el desarrollo. En particular, combinar resultados de igniciones independientes constituye una hipótesis de aproximación que requiere contraste, por lo que la coincidencia espacial entre escenarios no se presentará como una probabilidad calibrada de afectación.

La aplicación apoyará la interpretación del comandante o coordinador y permitirá explorar los escenarios como material de capacitación, sin asumir decisiones operativas automáticas ni incorporar un simulador táctico de colocación de recursos. El alcance tampoco incluye por ahora recepción de reportes comunitarios, gestión de recursos, integración de sensores o interoperabilidad institucional, aunque estas necesidades se conservarán como hallazgos de las entrevistas. La preparación anticipada reduce la necesidad de simulación durante el incidente, pero no acredita funcionamiento completo sin conexión, pues la consulta meteorológica y las condiciones de contingencia todavía requieren definición. El diseño anterior del póster~\cite{poster2026} sirve como antecedente documental, mientras que esta propuesta recoge las decisiones posteriores de escritorio, WFDS exclusivo y densidad espacial variable.
